\chapter{Запрет канонического автономного унитарного оператора часов из одного канала}
\label{ch:v8-canonical-autonomous-clock-unitary-extension-no-go-gate}

\section{Постановка сильнейшего внутреннего теста}

Предыдущая глава построила стационарный конечный родитель истории, но оставила
открытым полный унитарный такт. Проверим, снимают ли эту свободу уже
полученные минимальность среды, калибровочная ковариантность, вещественная структура и
чётность.

Пусть
\begin{equation}
 W:\mathbb C^{21}\longrightarrow\mathbb C^{21}\otimes\mathbb C^{13}
 =\mathbb C^{273}
\end{equation}
есть минимальная Стайнспринг-изометрия, а $P_W=WW^*$ --- проектор на её
образ. Любой полный унитарий $U_0$ с условием
\begin{equation}
 U_0(\psi\otimes|0\rangle)=W\psi
 \label{eq:v8-clock-unitary-extension-condition}
\end{equation}
реализует тот же одношаговый канал.

\section{Фазовое семейство дополнения}

Для $|z|=1$ определим
\begin{equation}
 V_z=P_W+z(I-P_W),
 \qquad U_z=V_zU_0.
 \label{eq:v8-clock-complement-phase-family}
\end{equation}
Тогда точно
\begin{equation}
 V_zW=W,
 \qquad
 U_z(\psi\otimes|0\rangle)=W\psi.
\end{equation}
Следовательно, все $U_z$ дают одну и ту же редуцированную Краус-карту.

Поскольку $W$ ковариантен, его образ является инвариантным подпространством и
$P_W$ коммутирует с полным представлением калибровочной группы. Поэтому $V_z$ также
ковариантен. Комплексная свобода содержит как минимум $U(1)$.

Вещественная и чётная типизация уменьшают её, но не уничтожают: значения
\begin{equation}
 z=+1,\qquad z=-1
\end{equation}
вещественны, чётны, ковариантны и дают разные полные унитарии при одном
канале. Полная произвольная часть живёт на дополнении размерности
$273-21=252$ и до дополнительных условий образует $U(252)$.

\section{Следствие для гамильтониана часов}

Если записывать такт как $U=e^{-i\delta t H}$, разные продолжения дают
разные полные гамильтонианы и разные логарифмические ветви, не меняя
наблюдаемый одношаговый канал на подготовленном вакуумном секторе среды.
Поэтому ни канал, ни его минимальный ранг Крауса, ни калибровочная ковариантность
и чётность вещественной структуры не выбирают автономный гамильтониан часов и длительность такта.

\begin{theorem}[Ковариантная неединственность полного унитария]
Минимальная калибровочно-ковариантная изометрия Стайнспринга Тома VIII допускает
непрерывное семейство полных унитарных продолжений, реализующих один и тот
же канал. Даже после вещественной и чётной типизации остаются по меньшей мере два
различных продолжения $z=\pm1$.
\end{theorem}

\begin{nogocriterion}[Запрет вывода автономных часов из Краус-карты]
Нельзя получать единственный гамильтониан часов Пейджа--Вуттерса простым
логарифмированием канала или произвольно выбранного Стайнспринг-унитария.
Требуется новый микроскопический гамильтониан взаимодействия, действие часов
или независимое физическое условие на дополнении $W\mathcal H$.
\end{nogocriterion}

\begin{protocol}
\begin{enumerate}
 \item системная размерность: \textbf{$21$};
 \item минимальная среда: \textbf{$13$};
 \item полный носитель одного такта: \textbf{$273$};
 \item образ изометрии: \textbf{$21$};
 \item свободное дополнение: \textbf{$252$};
 \item полное семейство продолжений: \textbf{$U(252)$};
 \item размерность полного семейства: \textbf{$63504$ вещественных измерений};
 \item ковариантное фазовое подсемейство: \textbf{$U(1)$};
 \item одинаковый редуцированный канал: \textbf{точно};
 \item вещественно-чётные различные представители: \textbf{$z=\pm1$};
 \item единственный автономный унитарий: \textbf{запрет};
 \item конечный условный мост истории: \textbf{сохранён};
 \item физическая длительность такта: \textbf{не выведена};
 \item следующий тест: \textbf{микроскопический гамильтониан взаимодействия
 или действие часов}.
\end{enumerate}
\end{protocol}

Машинный сертификат сохранён в
\path{s2t_v8_canonical_autonomous_clock_unitary_extension_no_go_gate_results.json}.

\begin{thebibliography}{9}
\bibitem{v8-clock-covariant-stinespring}
D.~Verdon,
\emph{A covariant Stinespring theorem},
Journal of Mathematical Physics 63 (2022), 091705; arXiv:2108.09872.
\bibitem{v8-clock-unitary-attal-pautrat}
S.~Attal, Y.~Pautrat,
\emph{From repeated to continuous quantum interactions},
Annales Henri Poincar\'e 7 (2006), 59--104.
\end{thebibliography}